% Part: sets-functions-relations
% Chapter: arithmetization
% Section: cuts
%
\documentclass[../../../include/open-logic-section]{subfiles}

\begin{document}

\olfileid{sfr}{arith}{cuts}
\olsection{Van $\Rat$ naar $\Real$}
De volledigheidseigenschap zegt eigenlijk dat elk punt $\alpha$ op de
getallenlijn die lijn netjes in twee helften verdeelt. De onderste
helft heeft $\alpha$ als kleinste bovengrens; de bovenste helft heeft
$\alpha$ als grootste ondergrens. Dedekind stelde voor om de reële
getallen uit de rationale getallen te \emph{maken}. Een rationaal getal
kun je schrijven als een breuk van gehele getallen, waarbij de noemer
niet nul is. Dedekind zag een
reëel getal als een \emph{snede} die de rationale getallen in tweeën
deelt. Zo stellen we $\sqrt{2}$ gelijk aan de \emph{snede} tussen de
rationale getallen $< \sqrt{2}$ en de rationale getallen $> \sqrt{2}$.
Dat kunnen we nog wat preciezer opschrijven. Als we de rationale getallen in twee
helften delen, is de \emph{onderste} helft genoeg om te weten welke
verdeling we hebben gemaakt. Daarom gebruiken we de volgende definitie:
\begin{defn}[Snede]
Een \emph{snede} $\alpha$ is een deel van de rationale getallen dat
niet leeg is en ook niet alle rationale getallen bevat. Als er een
getal in zit, zitten alle kleinere rationale getallen er ook in. Er is
geen grootste getal in de snede. Precies opgeschreven is $\alpha$ dus
een snede als:
\begin{enumerate}
	\item \emph{er iets in zit, maar niet alles}: $\emptyset \neq \alpha \subsetneq \Rat$
	\item \emph{alle kleinere getallen er ook in zitten}: voor alle $p,q \in \Rat$ geldt: als $p < q \in \alpha$, dan $p \in \alpha$
	\item \emph{er geen grootste getal is}: voor elke $p \in \alpha$ is er een $q \in \alpha$ waarvoor $p < q$
\end{enumerate}
De verzameling $\Real$ bestaat dan uit alle sneden.
\end{defn}
We kunnen nu dus zeggen dat $\sqrt{2} = \Setabs{p \in \Rat}{p^2 < 2\text{ of }p < 0}$. We moeten nog wel controleren dat dit echt een
\emph{snede} is. Dat doen we pas in \olref[check]{sec}.
Nu we weten wat onze getallen zijn, moeten we vastleggen hoe de
eenvoudige functies en relaties voor die getallen werken. We beginnen
met deze regel:
\begin{align*}
	\alpha \leq \beta \text{ precies dan als }\alpha \subseteq \beta
\end{align*}
Met deze regel voor de volgorde kunnen we het belangrijkste resultaat
\emph{formuleren}: de verzameling sneden heeft de
volledigheidseigenschap. Dat betekent hier het volgende. Als $S$ een
niet-lege verzameling sneden is
en een bovengrens heeft, dan heeft $S$ ook een kleinste bovengrens.
Volledig uitgeschreven: er is een snede $\lambda$ die een bovengrens
van $S$ is, oftewel\ $(\forall \alpha \in S)\alpha \subseteq
\lambda$. Bovendien is $\lambda$ de kleinste snede die dat doet,
oftewel\ $(\forall \beta \in \Real)((\forall \alpha \in S)\alpha \subseteq \beta \lif \lambda \subseteq \beta)$. Dit
is het bewijs:
\begin{thm}\ollabel{realcompleteness}
De verzameling sneden heeft de volledigheidseigenschap.
\end{thm}
\begin{proof}
Neem een willekeurige niet-lege verzameling sneden $S$ met een
bovengrens. Stel $\lambda = \bigcup S$.
%\Setabs{p \in \Rat}{(\exists \alpha \in S)p \in \alpha}$$
We laten eerst zien dat $\lambda$ een snede is:
\begin{enumerate}
\item $S$ is niet leeg. Daarom zit er minstens één snede in
$S$, zodat $\emptyset \neq \lambda$. Alles wat in $S$ zit, is een snede,
dus $\lambda \subseteq \Rat$. Omdat $S$ een bovengrens heeft, is er
een rationaal getal $p \in \Rat$ dat in geen enkele snede $\alpha \in S$
zit. Daarom
is $p\notin \lambda$ en dus $\lambda \subsetneq \Rat$.
\item Neem aan dat $p < q \in \lambda$. Dan is er een snede $\alpha \in S$
waarvoor $q \in \alpha$. Omdat $\alpha$ een snede is, geldt
$p \in \alpha$. Dus $p \in \lambda$.
\item Neem aan dat $p \in \lambda$. Dan is er een snede $\alpha \in S$
waarvoor $p \in \alpha$. Omdat $\alpha$ een snede is, is er een
$q \in \alpha$ waarvoor $p < q$. Dus $q \in \lambda$.
\end{enumerate}
Daarmee is de bewering bewezen. Verder is het duidelijk dat $(\forall \alpha \in S)\alpha
\subseteq \bigcup S = \lambda$, oftewel\ $\lambda$ is een bovengrens van $S$. Neem nu aan dat $\beta \in \mathbb{R}$ ook een bovengrens is, oftewel\ $(\forall \alpha \in S)\alpha \subseteq \beta$. Kies een willekeurige $p \in \Rat$. Als $p \in \lambda$, dan is er een $\alpha \in S$ met $p \in \alpha$, en dus geldt ook $p \in \beta$. Dit geldt voor elk rationaal getal, zodat $\lambda \subseteq \beta$. Daarom is $\lambda$ de \emph{kleinste} bovengrens van $S$.
\end{proof}
We hebben nu dus sneden die aan de volledigheidseigenschap voldoen.
Je kunt dat ook zo zeggen: er zitten geen “gaten” tussen onze sneden.
(Als we nog eens “sneden” zouden maken, maar nu in de reële getallen
in plaats van de rationale, zouden daar geen interessante nieuwe
dingen uit komen.)
Nu moeten we vastleggen hoe de bewerkingen op de reële getallen werken.
Eerst nemen we de rationale getallen op in de reële getallen. We
spreken af dat $p_\Real = \Setabs{q \in \Rat}{q < p}$ voor elke
$p \in \Rat$. Daarna definiëren we:
\begin{align*}
	\alpha + \beta &= \Setabs{p + q}{p \in \alpha \land q \in \beta}\\
%	\alpha - \beta &\defis \Setabs{p - q}{p \in \alpha \land q \in \Rat \setminus \beta}\\
	\alpha \times \beta &=
	\Setabs{p \times q}{0 \leq p \in \alpha \land 0 \leq q \in \beta} \cup 0_\Real & \text{als }\alpha, \beta \geq 0_\Real
%	\alpha \div \beta &\defis
%	\Setabs{\nicefrac{p}{q}}{p \in \alpha \land q \in \Rat \setminus \beta}  & \text{if }\alpha \geq 0_\Real, \beta > 0_\Real
\end{align*}
Verder spreken we voor elk reëel getal af dat nul maal dat getal nul
is en dat het getal maal nul ook nul is. Voor de andere gevallen van
vermenigvuldiging spreken we eerst het volgende af:
\begin{align*}
	-\alpha &= \Setabs{p - q}{p < 0 \land q \notin \alpha}
\end{align*}
Daarna spreken we af:
\begin{align*}
	\alpha \times \beta &\defis
	\begin{cases}
		\mathord{-}\alpha \times \mathord{-}\beta &\text{als }\alpha < 0_\Real\text{ en }\beta < 0_\Real\\
		\mathord{-}(\mathord{-}\alpha \times \beta) &\text{als }\alpha < 0_\Real \text{ en }\beta > 0_\Real\\
		\mathord{-}(\alpha \times \mathord{-}\beta) &\text{als }\alpha > 0_\Real \text{ en }\beta < 0_\Real
	\end{cases}
\end{align*}
We moeten nog controleren dat elk van deze regels altijd een snede
oplevert. Daarna moeten we laten zien dat de zo gemaakte sneden precies
doen wat ze horen te doen. Dat is niet moeilijk, maar het bewijs is wel
lang. Ook dit behandelen we in \olref[check]{sec}.
\end{document}